% GENERATED FILE. Edit canonical lessons, then run npm run generate:books.
% canonical-source-hash: fnv1a64:43be5688d7802f25
% canonical-lessons: ZH-C470-zisi, ZH-C470-guzhi, ZH-C470-jiaoao, ZH-C470-mahu, ZH-C470-ben

\chapter{Faults}
\label{ch:zh-faults}
\hlchaptermodality{470}

\begin{chapteropening}
\textbf{By the end of this chapter:} \emph{I can say someone is selfish, stubborn, a bit full of himself or sloppy, and that I feel stupid.}

Complete the last lesson of chapter 470: I can say someone is selfish, stubborn, a bit full of himself or sloppy, and that I feel stupid.
\end{chapteropening}

\section[zìsī]{zìsī — selfish}

\label{lesson:ZH-C470-zisi}



\begin{quote}
Before the new one: say the Mandarin for steady and reliable, then the Mandarin for an impatient sort.
\end{quote}

\subsection*{You'll want to know: zìsī}
\emph{zìsī} — \textquotedblleft{}selfish\textquotedblright{}.

\begin{quote}
\emph{wǒ zuì bù xǐhuan zìsī de rén} — Selfish people are the ones I like least.
\end{quote}

Its characters come later; for now it is a word for the ear.

\begin{cousinweb}
\emph{zìsī} is built from \zh{自} 'self' and \emph{sī} 'private', both inherited from Middle Chinese, \emph{dzijH} and \emph{sij} (Guangyun, in Baxter's transcription), which glosses the second 'not public'.
\end{cousinweb}

\begin{grammarlens}[title={What you've built}]
One more word for saying how it was.
\end{grammarlens}

\subsection*{Your turn}
\begin{itemize}
\raggedright
  \item \emph{Say it:} \emph{zìsī}
  \item \emph{Say it:} \emph{zìsī}, once more
  \item \emph{Say it:} \emph{wǒ zuì bù xǐhuan zìsī de rén} — Selfish people are the ones I like least.
  \item \emph{Recall:} say the Mandarin for steady and reliable, then the Mandarin for an impatient sort, then say \emph{zìsī} again
\end{itemize}

\begin{tcolorbox}[breakable,colback=teal!4,colframe=teal!35!black,title={Before you move on}]
What does zìsī mean? (\textquotedblleft{}Selfish\textquotedblright{}.) Say it once more.
\end{tcolorbox}
\section[gùzhí]{gùzhí — stubborn}

\label{lesson:ZH-C470-guzhi}



\begin{quote}
Before the new one: say the Mandarin for an impatient sort, then the Mandarin for selfish.
\end{quote}

\subsection*{You'll want to know: gùzhí}
\emph{gùzhí} — \textquotedblleft{}stubborn\textquotedblright{}.

\begin{quote}
\emph{wǒ yéye bǐ wǒ nǎinai gùzhí} — My grandfather is more stubborn than my grandmother.
\end{quote}

Its characters come later; for now it is a word for the ear.

\begin{cousinweb}
\emph{gùzhí} is built from \emph{gù} 'firm' and \emph{zhí} 'to hold', both inherited from Middle Chinese, \emph{kuH} and \emph{tsyip} (Guangyun, in Baxter's transcription), which glosses the first 'firm' and the second 'to hold' and 'to keep': holding on to one's view.
\end{cousinweb}

\begin{grammarlens}[title={What you've built}]
One more word for saying how it was.
\end{grammarlens}

\subsection*{Your turn}
\begin{itemize}
\raggedright
  \item \emph{Say it:} \emph{gùzhí}
  \item \emph{Say it:} \emph{gùzhí}, once more
  \item \emph{Say it:} \emph{wǒ yéye bǐ wǒ nǎinai gùzhí} — My grandfather is more stubborn than my grandmother.
  \item \emph{Recall:} say the Mandarin for an impatient sort, then the Mandarin for selfish, then say \emph{gùzhí} again
\end{itemize}

\begin{tcolorbox}[breakable,colback=teal!4,colframe=teal!35!black,title={Before you move on}]
What does gùzhí mean? (\textquotedblleft{}Stubborn\textquotedblright{}.) Say it once more.
\end{tcolorbox}
\section[jiāo'ào]{jiāo'ào — proud; full of oneself}

\label{lesson:ZH-C470-jiaoao}



\begin{quote}
Before the new one: say the Mandarin for selfish, then the Mandarin for stubborn.
\end{quote}

\subsection*{You'll want to know: jiāo'ào}
\emph{jiāo'ào} — \textquotedblleft{}proud; full of oneself\textquotedblright{}.

\begin{quote}
\emph{tā hěn cōngming kěshì yǒudiǎnr jiāo'ào} — He's very clever, but a bit full of himself.
\end{quote}

Its characters come later; for now it is a word for the ear.

\begin{cousinweb}
\emph{jiāo'ào} is built from \emph{jiāo} and \emph{ào} 'haughty', both inherited from Middle Chinese, \emph{kjew} and \emph{ngawH} (Guangyun, in Baxter's transcription). The dictionary glosses the first as a horse six feet tall and the second 'arrogant'; the sense 'proud' for the first is not given there.
\end{cousinweb}

\begin{grammarlens}[title={What you've built}]
One more word for saying how it was.
\end{grammarlens}

\subsection*{Your turn}
\begin{itemize}
\raggedright
  \item \emph{Say it:} \emph{jiāo'ào}
  \item \emph{Say it:} \emph{jiāo'ào}, once more
  \item \emph{Say it:} \emph{tā hěn cōngming kěshì yǒudiǎnr jiāo'ào} — He's very clever, but a bit full of himself.
  \item \emph{Recall:} say the Mandarin for selfish, then the Mandarin for stubborn, then say \emph{jiāo'ào} again
\end{itemize}

\begin{tcolorbox}[breakable,colback=teal!4,colframe=teal!35!black,title={Before you move on}]
What does jiāo'ào mean? (\textquotedblleft{}Proud; full of oneself\textquotedblright{}.) Say it once more.
\end{tcolorbox}
\section[mǎhu]{mǎhu — sloppy; slapdash}

\label{lesson:ZH-C470-mahu}



\begin{quote}
Before the new one: say the Mandarin for stubborn, then the Mandarin for full of oneself.
\end{quote}

\subsection*{You'll want to know: mǎhu}
\emph{mǎhu} — \textquotedblleft{}sloppy; slapdash\textquotedblright{}.

\begin{quote}
\emph{tā zuò zuòyè hěn mǎhu} — He's slapdash about his homework.
\end{quote}

Its characters come later; for now it is a word for the ear.

\begin{cousinweb}
\emph{mǎhu} is built from \zh{马} 'horse' and \emph{hǔ} 'tiger', both inherited from Middle Chinese, \emph{maeX} and \emph{xuX} (Guangyun, in Baxter's transcription). How a horse and a tiger came to mean 'careless' is not settled; the second syllable is said light.
\end{cousinweb}

\begin{grammarlens}[title={What you've built}]
One more word for saying how it was.
\end{grammarlens}

\subsection*{Your turn}
\begin{itemize}
\raggedright
  \item \emph{Say it:} \emph{mǎhu}
  \item \emph{Say it:} \emph{mǎhu}, once more
  \item \emph{Say it:} \emph{tā zuò zuòyè hěn mǎhu} — He's slapdash about his homework.
  \item \emph{Recall:} say the Mandarin for stubborn, then the Mandarin for full of oneself, then say \emph{mǎhu} again
\end{itemize}

\begin{tcolorbox}[breakable,colback=teal!4,colframe=teal!35!black,title={Before you move on}]
What does mǎhu mean? (\textquotedblleft{}Sloppy; slapdash\textquotedblright{}.) Say it once more.
\end{tcolorbox}
\section[bèn]{bèn — slow-witted; clumsy}

\label{lesson:ZH-C470-ben}



\begin{quote}
Before the new one: say the Mandarin for full of oneself, then the Mandarin for slapdash.
\end{quote}

\subsection*{You'll want to know: bèn}
\emph{bèn} — \textquotedblleft{}slow-witted; clumsy\textquotedblright{}.

\begin{quote}
\emph{wǒ gǎnjué zìjǐ hěn bèn} — I feel really stupid.
\end{quote}

Its characters come later; for now it is a word for the ear.

\begin{cousinweb}
\emph{bèn} is inherited from Middle Chinese \emph{bwonX} (Guangyun, in Baxter's transcription), glossed as the inside of a bamboo stem; the dictionary also records a fat man nicknamed with this word. The sense 'stupid' is not given there, so how it came is not settled.
\end{cousinweb}

\begin{grammarlens}[title={What you've built}]
One more word for saying how it was.
\end{grammarlens}

\subsection*{Your turn}
\begin{itemize}
\raggedright
  \item \emph{Say it:} \emph{bèn}
  \item \emph{Say it:} \emph{bèn}, once more
  \item \emph{Say it:} \emph{wǒ gǎnjué zìjǐ hěn bèn} — I feel really stupid.
  \item \emph{Recall:} say the Mandarin for full of oneself, then the Mandarin for slapdash, then say \emph{bèn} again
\end{itemize}

\begin{tcolorbox}[breakable,colback=teal!4,colframe=teal!35!black,title={Before you move on}]
What does bèn mean? (\textquotedblleft{}Slow-witted; clumsy\textquotedblright{}.) Say it once more.
\end{tcolorbox}
